\section*{Unit 1}
\addcontentsline{toc}{section}{Unit 1}
Empat berkas berikut digunakan dalam Unit 1. Setiap berkas tetap mengikuti lisensinya sendiri; tautan sumber dan lisensi dapat diklik.
\begin{description}
\item[\texttt{3d-function-6.svg}] MartinThoma (karya asli).
\href{https://commons.wikimedia.org/wiki/File:3d-function-6.svg}{Halaman sumber di Wikimedia Commons}; lisensi \href{https://creativecommons.org/licenses/by/3.0}{CC BY 3.0}.
\item[\texttt{Great circle passing through two points.svg}] HaEr48; karya turunan dari Polar angle to spherical side.svg oleh Episcophagus.
\href{https://commons.wikimedia.org/wiki/File:Great_circle_passing_through_two_points.svg}{Halaman sumber di Wikimedia Commons}; lisensi \href{https://creativecommons.org/licenses/by-sa/4.0}{CC BY-SA 4.0}.
\item[\texttt{2019-07-Helix.jpg}] Ag2gaeh (karya asli).
\href{https://commons.wikimedia.org/wiki/File:2019-07-Helix.jpg}{Halaman sumber di Wikimedia Commons}; lisensi \href{https://creativecommons.org/licenses/by-sa/4.0}{CC BY-SA 4.0}.
\item[\texttt{Planned flight map of the Oiseau Blanc.svg}] Pethrus; karya turunan dari BlankMap-World8.svg oleh AMK1211.
\href{https://commons.wikimedia.org/wiki/File:Planned_flight_map_of_the_Oiseau_Blanc.svg}{Halaman sumber di Wikimedia Commons}; lisensi \href{https://creativecommons.org/licenses/by-sa/3.0}{CC BY-SA 3.0}.
\end{description}
Identitas revisi, ukuran, SHA-1 Wikimedia, SHA-256, URL berkas asli, dan hash turunan cetak tercatat dalam manifest media yang disertakan.
